\clearpage
\section{LTspice double-pulse test}
The supplied EPCGaNLibrary archive contains a transistor subcircuit, not a
complete inverter module model. Two EPC2361 instances form a half-bridge:
the low-side DUT receives two pulses and the high-side gate is held off.
A 1 microhenry load inductor from DC positive to the switching node establishes
the current. First pulse duration is approximately $LI/V$; actual currents are
measured at first turn-off and second turn-on and stored separately.

The bench explicitly includes bus-side power-loop inductance, additional
common-source inductance, gate-loop inductance, loop resistance and a finite
driver edge with series resistance. It is a single-device-per-position test,
not an electrically scaled parallel-bank simulation. Temperature is fixed.
The completed run uses the following settings (resistances in ohms, parasitic
inductances in nH, edge times in ns):
\begin{center}\small\begin{tabular}{ll}
@DPTINPUTS@
\end{tabular}\end{center}
Internal model gate resistance is 0.4 ohm; external resistance and the assumed
driver resistance are additional. Separate ON/OFF driver impedances are not
modeled in this initial sweep.
\subsection{Energy measurement definitions}
The CSV reports two distinct quantities in microjoules:
\begin{itemize}
\item Signed drain-terminal energy: $\int v_{DS}i_Ddt$ over a fixed post-command
window. This includes a change in stored charge energy; individual $E_{on}$
or $E_{off}$ terminal integrals are not automatically junction heat.
\item Channel/access-resistance dissipation: the unchanged vendor electrical
model is observed with an extra, non-feedback output. It sums channel power
and drain/source resistor power. The turn-on value subtracts normal on-state
conduction over the window, using a separate DC sweep of the same model at
5 V gate drive and the same temperature, evaluated at the instantaneous load current.
Turn-off retains delayed channel conduction. Raw and baseline-subtracted values
are both saved. Gate resistance and gate leakage dissipation are excluded.
\end{itemize}
The default window is 150 ns. The paired first-off/second-on currents are close
but not exactly identical. This excess-energy convention is documented rather
than presented as a manufacturer's standardized $E_{on}/E_{off}$ specification.
The reverse-conducting complementary device is not included in the DUT heat
table; a complete inverter energy map must account for its actual gating and
deadtime. These tables are therefore not automatically substituted into the
analytical module model. The latter still uses its labeled assumptions.

\subsection{Slopes and voltage-margin screening}
Raw sample derivatives give maximum absolute $dv/dt$ and $di/dt$ for each edge.
The CSV also reports maximum slopes over a 1 ns interval to reduce sensitivity
to small numerical steps. These are bandwidth-dependent numbers, not universal
device constants. Actual waveforms and time-step sensitivity are retained.

Peak $V_{DS}$ is checked on both devices against a 90 V design ceiling and
the 100 V continuous absolute rating. Both package gate-source voltages are
checked against -4 V and +6 V. $\Delta V\sim L\,di/dt$ is an intuition for the
inductive contribution, not a replacement for the complete ringing waveform.
The swept bus inductance is additional to 0.2 nH common-source inductance.
No safety conclusion is inferred from simulation survival: the model has no
destructive breakdown or avalanche qualification mechanism. Voltage margin
must be rechecked with PCB-extracted parasitics, driver behavior and hardware.

@DPTRESULT@
\begin{center}\includegraphics[width=\linewidth]{../../results/double_pulse/nominal_waveforms.pdf}\end{center}
\clearpage
\begin{center}\includegraphics[width=\linewidth]{../../results/double_pulse/sensitivity.pdf}\end{center}
\subsection{Limitations of the supplied model}
The library history introduces EPC2361 in March 2024. Its typical fit is not
automatically a match to every July 2026 datasheet curve or production corner.
The subcircuit includes nonlinear charge and temperature-dependent electrical
parameters, but no package/PCB inductors, thermal network, explicit avalanche
breakdown or dynamic trapping model. Charge elements do not establish real
Coss hysteresis loss. Simulated slopes, energies and overshoot depend strongly
on added parasitics and driver assumptions. Numerical convergence at the
nominal point does not validate the physical model or every sweep corner.
